% paper/sections/04f_face_jet.tex
%
% §4.7  面ジェット $\FaceJet{u}$：基底 CCD 出力の面 API（FCCD と等価）
%        本節では数学的定義のみを与える；消費先章での具体使用は章末まとめ §\ref{sec:ccd_summary} を参照．
%
% ═══════════════════════════════════════════════════════════════════════════════
\subsection{\texorpdfstring{面ジェット $\FaceJet{u}$：基底 CCD 出力の面 API（FCCD と等価）}{面ジェット J\_f(u)：基底 CCD 出力の面 API（FCCD と等価）}}
\label{sec:face_jet_def}
% ═══════════════════════════════════════════════════════════════════════════════

本節では，基底 CCD（§\ref{sec:ccd_def}）が返す節点 Hermite ペア
$(u_i,\,q_i = u''_i)$ から面値・面勾配・面二次導関数を構成する
\textbf{3 成分 API}：
\[
  \FaceJet{u} \;=\; (u_f,\,u'_f,\,u''_f)
\]
を\textbf{面ジェット}（face jet）として定義する．
構成式は §\ref{sec:fccd_def} の FCCD 公式と数式上\textbf{等価}であり，
FCCD は本 API を面中心スキームとして直接実装した形態と見なせる；
UCCD6（§\ref{sec:uccd6_def}）など他の節点中心 CCD 派生も
本 API を経由して面位置の表現を得る．
本節では数学的定義と各成分の数値的意味のみを与え，
後続章（§\ref{sec:time_int} 以降）では本 API を
\textbf{入力契約}（input contract）として再利用する；
個別の用途と機構は §\ref{sec:ccd_summary} の比較表に集約する．

% -------------------------------------------------------------------------------
\subsubsection{定義：面ジェット演算子}
\label{sec:face_jet_formula}
% -------------------------------------------------------------------------------

節点上の値 $u_i$ と CCD が返す節点 Hermite ペア $(u'_i, u''_i = q_i)$ を入力とし，
面 $x_{i-1/2}$ 上の 3 組を：
\begin{equation}
  \FaceJet{u}_{i-1/2}
  \;=\;
  \bigl(P_f u,\; G_f u,\; Q_f u\bigr)_{i-1/2}
  \;\equiv\;
  \bigl(u_{i-1/2},\; u'_{i-1/2},\; u''_{i-1/2}\bigr)
  \label{eq:face_jet_def}
\end{equation}
で定義する．ここで $P_f, G_f, Q_f$ はそれぞれ面値・面勾配・面二次導関数を
コンパクト Hermite 公式：
\begin{align}
  P_{i-1/2}\,u &=
  \frac{u_{i-1}+u_i}{2}
  -\frac{H_i^2}{16}(q_{i-1}+q_i), \label{eq:face_value_formula}\\
  G_{i-1/2}\,u &=
  \frac{u_i-u_{i-1}}{H_i}
  -\frac{H_i}{24}(q_i-q_{i-1}), \label{eq:face_grad_formula}\\
  Q_{i-1/2}\,u &=
  \tfrac{1}{2}(q_{i-1}+q_i) \label{eq:face_second_formula}
\end{align}
で構成する（$q_i = u''_i$ は CCD の出力である二次導関数）．
式~\eqref{eq:face_value_formula}–\eqref{eq:face_grad_formula} は $\Ord{H^4}$，
式~\eqref{eq:face_second_formula} は $\Ord{H^2}$ の精度を持つ；
二次導関数成分が意図的に低次で構成されている理由は，
これが独立未知数ではなく診断および接続条件評価に用いる\textbf{補助量}だからである．
\textbf{連立構造（coupling）：}
$P_{i-1/2}$, $G_{i-1/2}$ は $u_{i-1}, u_i$ と $q_{i-1}, q_i$ の双方を参照するため，
面ジェットの評価は CCD ブロック Thomas（§\ref{sec:ccd_def}）の出力 $\{q_i\}$ に
動的に依存する；面値・面勾配・面二次の 3 成分を独立に評価することはできず，
CCD 出力との連立が必須である．

% -------------------------------------------------------------------------------
\subsubsection{3 成分の数値的意味}
\label{sec:face_jet_roles}
% -------------------------------------------------------------------------------

面ジェット 3 成分の数値的意味（精度オーダーと幾何配置）は以下の通り：

\begin{center}
\begin{tabular}{l p{0.40\linewidth} l}
\hline
\textbf{成分} & \textbf{数値的意味} & \textbf{精度} \\ \hline
$u_f$    & コンパクト Hermite 公式による面値        & $\Ord{H^4}$ \\
$u'_f$   & コンパクト Hermite 公式による面勾配       & $\Ord{H^4}$ \\
$u''_f$  & 節点二次導関数の単純面平均（補助量）        & $\Ord{H^2}$ \\
\hline
\end{tabular}
\end{center}

\noindent
3 成分が単一の面上で\textbf{同一の演算子族}から構成されることが
本 API の核心である．成分ごとに別ステンシル・別補間を用いると
微分演算子の離散的随伴性 $\langle D_h u, v\rangle = -\langle u, D_h^{*} v\rangle$ が
破れ，下流の連立離散において自己整合性を失う；
面ジェットはこの整合性を\textbf{API レベル}で保証する．

% -------------------------------------------------------------------------------
\subsubsection{API としての設計帰結}
\label{sec:face_jet_design_consequences}
% -------------------------------------------------------------------------------

\begin{itemize}
  \item \textbf{面ジェットは公開契約（public contract）．}
        現実装は既存の CCD ノード解 $(u'_i, u''_i)$ を再利用する\textbf{加算的}
        （additive）構成であるが，将来的にブロックシステムで
        $(u_f, u'_f, u''_f)$ を独立未知数に昇格しても
        呼び出し側インタフェースは不変である．
  \item \textbf{不連続生フィールドに面値原語を直接適用しない．}
        ジャンプ処理なしの $u_f$ 適用は，
        $H^2 q$ 項（式~\eqref{eq:face_value_formula} 第 2 項）が
        不連続を Gibbs 振動として輸送する；
        本 API の数値的妥当性は被作用関数の $C^k$ 滑らかさに依存する．
  \item \textbf{連立構造の局所性．}
        ブロック Thomas（§\ref{sec:ccd_matrix}）が返す節点解 $\{q_i\}$ から
        各面が独立に評価されるため，計算量は $\Ord{N}$ で
        基底 CCD と同オーダーである．
\end{itemize}

\phantomsection\label{sec:face_jet_hfe}%
\phantomsection\label{sec:face_jet_ppe}%
\phantomsection\label{eq:hfe_upwind_plus}%
\phantomsection\label{eq:hfe_upwind_minus}%
\phantomsection\label{eq:face_ppe_flux}%

% ═══════════════════════════════════════════════════════════════════════════════
\subsection{CCD スキーム族の数値特性と略語：章末整理}
\label{sec:ccd_summary}
% ═══════════════════════════════════════════════════════════════════════════════

本章では，3 点ステンシル制約下で $\Ord{h^6}$ 精度を実現する
コンパクト差分スキーム族 — CCD・DCCD・FCCD・UCCD6 — と
その出力を統一する 3 成分 API である面ジェット $\FaceJet{u}=(u_f,u'_f,u''_f)$ を導出した．
本節では数値特性を比較表で整理し，後続章（§\ref{sec:time_int} 以降）への
参照を章名と節番号のみで示す．

% -------------------------------------------------------------------------------
\subsubsection*{数値特性比較表}
% -------------------------------------------------------------------------------

\begin{center}
{\footnotesize
\begin{tabular}{l l l l l l}
\hline
\textbf{スキーム} & \textbf{配置} & \textbf{散逸チャンネル} & \textbf{微分精度}$^\dagger$ & \textbf{安定性} & \textbf{計算量} \\ \hline
{\CCD} (§\ref{sec:ccd_def})              & 節点      & なし                          & $\Ord{h^6}$           & 中心差分（散逸ゼロ）         & $\Ord{N}$ \\
{\DCCD} (§\ref{sec:dccd_derivation})     & 節点      & 外殻 post-filter              & $\Ord{h^6}$$^\ddagger$ & spectral 非増幅（係数拘束下）   & $\Ord{N}$ \\
{\FCCD} (§\ref{sec:fccd_def})            & \textbf{面}$^\flat$ & なし（散逸ゼロ）          & $\Ord{H^4}$           & 基底 CCD 由来               & $\Ord{N}$ \\
{\UCCD[6]} (§\ref{sec:uccd6_def})        & 節点      & \textbf{演算子内部}（hyperviscosity） & $\Ord{h^6}$  & CN で無条件安定（周期格子）    & $\Ord{N}$ \\
$\FaceJet{u}$ (§\ref{sec:face_jet_def})  & 面（API） & ---                           & $\Ord{H^4}$/$\Ord{H^2}$$^\S$ & 基底スキーム継承   & $\Ord{N}$ post \\
\hline
\end{tabular}}
\end{center}

\noindent
{\footnotesize
$^\dagger$ 微分演算子の formal 精度．
$^\ddagger$ DCCD は微分演算子そのものは $\Ord{h^6}$ を保つが，
post-filter による移流フィルタ誤差 $\Ord{\varepsilon_d^{\mathrm{adv}} h^2}$ が別途付随し
（§\ref{sec:dccd_filter_theory}），移流の実効精度は $\Ord{h^2}$ に律速される．
$^\flat$ FCCD は面位置に値を置く面中心スキームだが Fourier シンボルは純虚数で散逸ゼロ；
散逸が必要な場面では速度向き数値流束または別フィルタを併用する．
$^\S$ 面ジェットは $u_f, u'_f$ が $\Ord{H^4}$，$u''_f$ が $\Ord{H^2}$．
}

% -------------------------------------------------------------------------------
\subsubsection*{後続章への参照一覧}
% -------------------------------------------------------------------------------

本章で導入した演算子・API は以下の章で再利用される（章名と節番号のみ；
個別機構の詳細は各章で展開）：

\begin{itemize}\setlength{\itemsep}{0pt}
  \item \S\ref{sec:time_int}（時間積分）：UCCD6 を Crank--Nicolson と組み合わせる．
  \item \S\ref{sec:grid}（非一様格子）：CCD/FCCD の非一様格子拡張．
  \item \S\ref{sec:advection}（移流）：CLS 移流に FCCD（面ジェット保存型），運動量移流に FCCD/UCCD6 を適用．
  \item \S\ref{sec:collocate}（圧力・速度連成）：FCCD 面ジェットを面演算子として再利用．
  \item \S\ref{sec:ppe_solve}（圧力 Poisson）：CCD 系微分作用素を PPE の離散化基底に用いる．
\end{itemize}

% -------------------------------------------------------------------------------
\subsubsection*{略語一覧}
% -------------------------------------------------------------------------------

\begin{description}\setlength{\itemsep}{0pt}
  \item[CCD] Combined Compact Difference（Chu--Fan 6 次基底）（§\ref{sec:ccd_def}）
  \item[DCCD] Dissipative CCD（外殻 spectral post-filter 付加）（§\ref{sec:dccd_derivation}）
  \item[FCCD] Face-centered CCD（面位置の値・導関数を $\Ord{H^4}$ で構成）（§\ref{sec:fccd_def}）
  \item[UCCD6] Upwind CCD-6（次数保存ハイパー粘性内蔵）（§\ref{sec:uccd6_def}）
  \item[面ジェット $\FaceJet{u}$] $(u_f, u'_f, u''_f)$ を一括して提供する 3 成分 API（§\ref{sec:face_jet_def}）
\end{description}

% ═══════════════════════════════════════════════════════════════════════════════
